\documentclass[12pt,a4paper,oneside]{ctexbook}
\usepackage[utf8]{inputenc}
\usepackage{geometry}
\geometry{left=2.5cm,right=2.5cm,top=3cm,bottom=3cm}
\usepackage{amsmath,amssymb}
\usepackage{hyperref}
\usepackage{booktabs}
\usepackage{longtable}
\usepackage{fancyhdr}
\usepackage{setspace}
\usepackage{enumitem}
\usepackage{listings}
\usepackage{xcolor}
\usepackage{float}

\lstset{
    basicstyle=\small\ttfamily,
    numbers=left,
    numberstyle=\tiny,
    frame=single,
    breaklines=true,
    backgroundcolor=\color{gray!5},
    tabsize=4,
    language=C++
}

\setlength{\parindent}{2em}
\setstretch{1.3}

\pagestyle{fancy}
\fancyhf{}
\fancyhead[LE,RO]{\leftmark}
\fancyfoot[C]{\thepage}

\hypersetup{colorlinks=true, linkcolor=black}

\begin{document}

\title{Modular Yielding Graphical Operators\\MYGO-DSP\\数字信号处理与滤波器设计平台}
\author{}
\date{版本 2.0 -- 2026年5月}
\maketitle
\thispagestyle{empty}

\frontmatter
\tableofcontents

\mainmatter

\chapter{项目概述}

\section{什么是 MYGO-DSP?}

MYGO-DSP 是一个企业级数字信号处理与滤波器设计平台, 基于 C++17 实现高性能 DSP 核心引擎, 通过 Electron + React 构建跨平台桌面应用界面。

M.Y.G.O 代表 Modular Yielding Graphical Operators（模块化生成图形算子）。核心理念是将信号处理单元（波形发生器、FIR/IIR 滤波器）抽象为统一的"算子"接口, 通过图形界面组合和实时控制。

\section{数字信号处理基础}

数字信号处理是将连续信号离散化为一系列采样点进行分析和处理的技术。一个正弦波离散信号可表示为:

\[
x[n] = A \cdot \sin(2\pi f n / f_s + \phi)
\]

其中 \(A\) 为幅度, \(f\) 为频率 (Hz), \(f_s\) 为采样率, \(\phi\) 为初始相位。

奈奎斯特采样定理指出:

\[
f_s \geq 2 f_{\max}
\]

采样率必须至少为目标信号最高频率的两倍才能无失真重建。

傅里叶变换是 DSP 最核心的数学工具, 将时域信号变换到频域:

\[
X[k] = \sum_{n=0}^{N-1} x[n] \cdot e^{-j2\pi kn/N}
\]

快速傅里叶变换 (FFT) 将计算复杂度从 \(O(N^2)\) 降至 \(O(N\log N)\)。

\section{系统架构}

系统采用三层分离架构:

\begin{itemize}
    \item \textbf{C++ 引擎层}: 信号处理核心。7 种波形生成、FIR/IIR 滤波、智能滤波器设计、FFT 分析。通过 stdin/stdout JSON 协议与前端通信。
    \item \textbf{Electron 桥接层}: 启动管理 DSP 子进程, 维护 IPC 命令队列, 转发请求与响应。
    \item \textbf{React 前端层}: macOS Big Sur 风格 UI。参数控制面板、ECharts 图表（时域、频谱、零极点）、智能设计弹窗。
\end{itemize}

\chapter{环境搭建}

\section{必要工具}

\begin{itemize}
    \item CMake $\geq$ 3.16: 构建系统生成器
    \item MinGW/GCC $\geq$ 8.1: C++17 编译器
    \item Node.js $\geq$ 18: JavaScript 运行时
    \item Visual Studio Code (推荐): 代码编辑器
    \item Git: 版本管理
\end{itemize}

\section{安装 CMake}

访问 \url{https://cmake.org/download/} 下载 Windows 安装包。安装时勾选"Add CMake to system PATH"。验证:

\begin{lstlisting}
cmake --version
\end{lstlisting}

\section{安装 MinGW/GCC}

下载 MinGW-w64, 或使用 Qt 自带的版本 (如 \texttt{D:\textbackslash Qt\textbackslash Tools\textbackslash mingw1120\_64})。将 \texttt{bin} 目录加入 PATH:

\begin{lstlisting}
g++ --version
\end{lstlisting}

需支持 C++17 标准 (GCC $\geq$ 8.1)。

\section{安装 Node.js}

访问 \url{https://nodejs.org/} 下载 LTS 版本:

\begin{lstlisting}
node --version
npm --version
\end{lstlisting}

\chapter{代码结构}

每个文件的作用:

\begin{itemize}
    \item \textbf{DSPMath.h}: 数学常量 \(\pi\), \(2\pi\)。避免不同编译器中 \texttt{M\_PI} 定义不一致的问题。
    \item \textbf{Buffer.h}: 环形缓冲区模板类, 用于数据流 FIFO 管理。线程安全（\texttt{std::mutex} 保护）。
    \item \textbf{IOperator.h}: 算子抽象接口（虚基类）。定义 \texttt{process()}, \texttt{reset()}, \texttt{setBypass()} 通用接口。波形发生器和滤波器都继承自此类。
    \item \textbf{WaveformGenerator}: 相位累加器驱动。支持 7 种波形, 使用 Mersenne Twister MT19937-64 生成噪声。
    \item \textbf{FIRFilter}: 窗函数法设计 + 直接型卷积实现。5 种窗函数, 4 种类型 (LP/HP/BP/BS)。
    \item \textbf{IIRFilter}: 模拟原型法设计 + 双线性变换。Biquad 级联实现, 直接 II 型转置结构。
    \item \textbf{AutoFilterDesigner}: 智能原型选择和阶数估算。椭圆→切比雪夫→巴特沃斯→FIR 凯泽窗回退。
    \item \textbf{Engine}: 单例调度器。管理所有波形/滤波器实例, 处理管线按"波形叠加→滤波器链"顺序执行。
    \item \textbf{main.cpp}: IPC 主循环。从 stdin 读取 JSON 命令, 解析后调用 Engine, 将结果 JSON 写入 stdout。
\end{itemize}

\chapter{C++ DSP 核心算法}

\section{波形发生器}

相位累加器 \(\phi[n]\) 以采样率递增:

\[
\phi[n] = \phi[n-1] + \frac{f}{f_s}
\]

各波形公式:

\begin{align*}
\text{正弦:}&\quad x[n] = A \sin(2\pi \phi[n] + \theta) \\
\text{方波:}&\quad x[n] = \begin{cases} A, & \text{frac}(\phi+\theta/2\pi)<D \\ -A, & \text{else} \end{cases} \\
\text{三角:}&\quad x[n] = 4A|\text{frac}(\phi+\theta/2\pi)-0.5| - A \\
\text{锯齿:}&\quad x[n] = 2A\cdot\text{frac}(\phi+\theta/2\pi) - A
\end{align*}

方波的傅里叶级数含丰富奇次谐波:

\[
x_{\text{square}}(t) = \frac{4A}{\pi}\sum_{k=1,3,5\ldots}\frac{\sin(2\pi k f t)}{k}
\]

\section{FIR 滤波器}

输出为输入与脉冲响应的卷积:

\[
y[n] = \sum_{k=0}^{N-1} h[k] x[n-k]
\]

窗函数法设计步骤:

\begin{enumerate}
\item 理想脉冲响应: \(h_{\text{ideal}}[n] = \sin(\omega_c(n-M))/\pi(n-M)\)
\item 加窗截断: \(h[n] = h_{\text{ideal}}[n] \cdot w[n]\)
\item 归一化: \(\sum h[n] = 1\)
\end{enumerate}

窗函数比较:

\[
\begin{array}{l|l|l}
\text{窗类型} & \text{主瓣宽度} & \text{旁瓣衰减} \\ \hline
\text{矩形} & 4\pi/N & -13\text{dB} \\
\text{汉明} & 8\pi/N & -43\text{dB} \\
\text{汉宁} & 8\pi/N & -32\text{dB} \\
\text{布莱克曼} & 12\pi/N & -58\text{dB}
\end{array}
\]

凯泽窗: \(w[n] = I_0(\beta\sqrt{1-(2n/(N-1)-1)^2})/I_0(\beta)\), 通过 \(\beta\) 调节主瓣/旁瓣权衡。

\section{IIR 滤波器}

差分方程:

\[
y[n] = \sum b_k x[n-k] - \sum a_k y[n-k]
\]

采用 Biquad 级联维持数值稳定性:

\[
H(z) = \prod_{k=1}^{K} \frac{b_{k0}+b_{k1}z^{-1}+b_{k2}z^{-2}}{1+a_{k1}z^{-1}+a_{k2}z^{-2}}
\]

模拟原型设计从 s 平面极点分布开始, 通过双线性变换 \(s = (2/T)(1-z^{-1})/(1+z^{-1})\) 映射到 z 平面。

\section{智能设计器}

阶数估算公式:

巴特沃斯:
\[
N = \left\lceil\frac{\log_{10}((10^{A_s/10}-1)/(10^{A_p/10}-1))}{2\log_{10}(\omega_s/\omega_p)}\right\rceil
\]

椭圆:
\[
N = \left\lceil\frac{K(k)K(\sqrt{1-k_1^2})}{K(\sqrt{1-k^2})K(k_1)}\right\rceil
\]

FIR 凯泽窗回退:
\[
N = \left\lceil\frac{A_s-7.95}{14.36\cdot\Delta f}\right\rceil,\quad \beta = \begin{cases}
0.1102(A_s-8.7), & A_s>50 \\
0.5842(A_s-21)^{0.4}+0.07886(A_s-21), & 21\leq A_s\leq 50
\end{cases}
\]

\chapter{IPC 通信协议}

\section{格式}

每行一个 JSON 对象:

\begin{lstlisting}
{"id":1,"cmd":"addWave","type":0,"frequency":440,"amplitude":0.5}
{"id":1,"status":"ok","data":0}
\end{lstlisting}

命令列表:

\begin{longtable}{p{3cm}p{3cm}p{7cm}}
\caption{命令一览}
\endfirsthead
\hline
\textbf{命令} & \textbf{参数} & \textbf{说明} \\
\hline
\endhead
\hline
\endfoot
init & sampleRate, bufferSize & 初始化 \\
addWave & type, frequency, amplitude, phase, dutyCycle & 加波形 (返回ID) \\
removeWave & wid & 删波形 \\
updateWave & wid, type, ... & 改波形参数 \\
addFilterIIR & prototype, passType, order, cutoff, ... & 加 IIR 滤波器 \\
addFilterFIR & order, window, cutoff, passType, ... & 加 FIR 滤波器 \\
removeFilter & fid & 删滤波器 \\
updateFilter & fid, ..., isIIR & 改滤波器参数 \\
setBypass & bid, bypass, target & 设置旁路 \\
process & (无) & 处理一帧 (128 采样点) \\
autoDesign & type, fs, fpass, fstop, Ap, As & 智能设计 \\
getPoleZero & fid & 获取零极点 \\
exit & (无) & 退出 \\
\hline
\end{longtable}

\chapter{构建}

\section{DSP 引擎}

\begin{lstlisting}
cd dsp-core
mkdir build && cd build
cmake .. -G "MinGW Makefiles"
mingw32-make -j4
\end{lstlisting}

静态链接: \texttt{-static -static-libgcc -static-libstdc++}。

\section{前端}

\begin{lstlisting}
cd frontend
npm install
npx vite build
\end{lstlisting}

\section{Electron 打包}

\begin{lstlisting}
npx --package @electron/asar asar pack dist+electron app.asar
\end{lstlisting}

将 \texttt{app.asar} 和 \texttt{mygo\_dsp\_engine.exe} 放入 Electron \texttt{resources} 目录。

\chapter{测试}

\section{单元测试 (Google Test)}

\begin{longtable}{p{3.5cm}p{1.5cm}p{8cm}}
\caption{38 个测试用例}
\endfirsthead
\hline
\textbf{模块} & \textbf{数量} & \textbf{覆盖} \\
\hline
\endhead
\hline
\endfoot
WaveformGenerator & 6 & 正弦输出, 频率精度, bypass, 占空比, 噪声, 重置 \\
FIRFilter & 5 & 直流响应, 脉冲响应, bypass, Kaiser \\
IIRFilter & 6 & 稳定性, 衰减, bypass, RBJ增益, 重置 \\
Engine & 5 & 单例, 初始化, 滤波链, bypass, 多波形 \\
AutoDesigner & 8 & LP/HP/BP/BS, FIR回退, 无效参数, 阶数估算 \\
\hline
\end{longtable}

\section{性能基准}

Intel i7-10750H, MinGW -O2:

\begin{longtable}{p{5cm}p{3cm}p{3cm}}
\caption{性能测试}
\endfirsthead
\hline
\textbf{操作} & \textbf{吞吐量} & \textbf{延迟} \\
\hline
\endhead
\hline
\endfoot
波形生成 & 18.9M 样本/秒 & 52.83 ns \\
FIR order=32 & 5.2M 样本/秒 & 193.25 ns \\
FIR order=64 & 3.1M 样本/秒 & 319.64 ns \\
FIR order=128 & 1.7M 样本/秒 & 599.30 ns \\
IIR order=4 & 45.4M 样本/秒 & 22.02 ns \\
IIR order=8 & 21.5M 样本/秒 & 46.56 ns \\
IIR order=12 & 13.5M 样本/秒 & 73.88 ns \\
完整管线 & 23928 帧/秒 & 41.79 $\mu$s \\
\hline
\end{longtable}

\chapter{理论验证}

\section{低通滤波}

250Hz 正弦 + 2000Hz 方波, 通过 500Hz 4 阶椭圆低通。

滤波后输出光滑正弦 (前 10 采样: \(4.7\times10^{-7},4.3\times10^{-6},\ldots\))。未滤波时含方波跳变 (\(0.748,0.756,0.163,\ldots\))。结论: 频率选择性正确。✅

\section{Bypass 功能}

滤波状态 (能量 42.96) 和 bypass 状态 (29.42) 显著不同。切换正常。✅

\section{FIR 稳定性}

全部输出为有限值, 无 NaN/Inf, 不发散。✅

\section{椭圆阶数估算}

fpass=500Hz, fstop=1000Hz, Ap=1dB, As=40dB 指标下自动选 order=4 椭圆。符合理论。✅

\chapter{零极点分析}

\section{IIR 极点}

模拟原型极点经双线性变换映射到 z 平面。所有极点必须位于单位圆内才稳定。Butterworth 全极点 (无有限零点), Chebyshev II 和 Elliptic 同时含极点和零点。

\section{FIR 零点}

FIR 滤波器是全零点滤波器。零点通过 Newton-Raphson 法求系数多项式根获得。

\section{实现}

\lstinline|getPoleZero| 命令返回所有滤波器的零极点。前端遍历 filters 数组, 合并显示。

\chapter{项目总结}

\section{技术指标}

\begin{itemize}
    \item DSP 引擎: 约 2500 行 C++17
    \item 前端: 约 1800 行 TypeScript + React
    \item 7 种波形, FIR 阶数 1-200, IIR 阶数 1-20
    \item 缓冲区 128 采样点
    \item 最大吞吐 45M 样本/秒 (IIR order=4)
\end{itemize}

\section{面试价值}

\begin{itemize}
    \item C++17 现代语法: 智能指针, 移动语义, Lambda, 模板
    \item 设计模式: 单例 (Engine), 策略 (IOperator 多态)
    \item 算法: 环形缓冲区, FFT, 多项式求根, 椭圆积分
    \item DSP 理论: z 变换, 双线性变换, 窗函数法, 滤波器设计
    \item 系统设计: 跨语言 IPC, 多进程架构, 实时管线
    \item 工程: CMake, Google Test, 性能基准
\end{itemize}

\end{document}
